% Options for packages loaded elsewhere
\PassOptionsToPackage{unicode}{hyperref}
\PassOptionsToPackage{hyphens}{url}
\documentclass[
  12pt,
]{article}
\usepackage{xcolor}
\usepackage[margin=2.5cm]{geometry}
\usepackage{amsmath,amssymb}
\setcounter{secnumdepth}{5}
\usepackage{iftex}
\ifPDFTeX
  \usepackage[T1]{fontenc}
  \usepackage[utf8]{inputenc}
  \usepackage{textcomp} % provide euro and other symbols
\else % if luatex or xetex
  \usepackage{unicode-math} % this also loads fontspec
  \defaultfontfeatures{Scale=MatchLowercase}
  \defaultfontfeatures[\rmfamily]{Ligatures=TeX,Scale=1}
\fi
\usepackage{lmodern}
\ifPDFTeX\else
  % xetex/luatex font selection
\fi
% Use upquote if available, for straight quotes in verbatim environments
\IfFileExists{upquote.sty}{\usepackage{upquote}}{}
\IfFileExists{microtype.sty}{% use microtype if available
  \usepackage[]{microtype}
  \UseMicrotypeSet[protrusion]{basicmath} % disable protrusion for tt fonts
}{}
\makeatletter
\@ifundefined{KOMAClassName}{% if non-KOMA class
  \IfFileExists{parskip.sty}{%
    \usepackage{parskip}
  }{% else
    \setlength{\parindent}{0pt}
    \setlength{\parskip}{6pt plus 2pt minus 1pt}}
}{% if KOMA class
  \KOMAoptions{parskip=half}}
\makeatother
\usepackage{color}
\usepackage{fancyvrb}
\newcommand{\VerbBar}{|}
\newcommand{\VERB}{\Verb[commandchars=\\\{\}]}
\DefineVerbatimEnvironment{Highlighting}{Verbatim}{commandchars=\\\{\}}
% Add ',fontsize=\small' for more characters per line
\usepackage{framed}
\definecolor{shadecolor}{RGB}{248,248,248}
\newenvironment{Shaded}{\begin{snugshade}}{\end{snugshade}}
\newcommand{\AlertTok}[1]{\textcolor[rgb]{0.94,0.16,0.16}{#1}}
\newcommand{\AnnotationTok}[1]{\textcolor[rgb]{0.56,0.35,0.01}{\textbf{\textit{#1}}}}
\newcommand{\AttributeTok}[1]{\textcolor[rgb]{0.13,0.29,0.53}{#1}}
\newcommand{\BaseNTok}[1]{\textcolor[rgb]{0.00,0.00,0.81}{#1}}
\newcommand{\BuiltInTok}[1]{#1}
\newcommand{\CharTok}[1]{\textcolor[rgb]{0.31,0.60,0.02}{#1}}
\newcommand{\CommentTok}[1]{\textcolor[rgb]{0.56,0.35,0.01}{\textit{#1}}}
\newcommand{\CommentVarTok}[1]{\textcolor[rgb]{0.56,0.35,0.01}{\textbf{\textit{#1}}}}
\newcommand{\ConstantTok}[1]{\textcolor[rgb]{0.56,0.35,0.01}{#1}}
\newcommand{\ControlFlowTok}[1]{\textcolor[rgb]{0.13,0.29,0.53}{\textbf{#1}}}
\newcommand{\DataTypeTok}[1]{\textcolor[rgb]{0.13,0.29,0.53}{#1}}
\newcommand{\DecValTok}[1]{\textcolor[rgb]{0.00,0.00,0.81}{#1}}
\newcommand{\DocumentationTok}[1]{\textcolor[rgb]{0.56,0.35,0.01}{\textbf{\textit{#1}}}}
\newcommand{\ErrorTok}[1]{\textcolor[rgb]{0.64,0.00,0.00}{\textbf{#1}}}
\newcommand{\ExtensionTok}[1]{#1}
\newcommand{\FloatTok}[1]{\textcolor[rgb]{0.00,0.00,0.81}{#1}}
\newcommand{\FunctionTok}[1]{\textcolor[rgb]{0.13,0.29,0.53}{\textbf{#1}}}
\newcommand{\ImportTok}[1]{#1}
\newcommand{\InformationTok}[1]{\textcolor[rgb]{0.56,0.35,0.01}{\textbf{\textit{#1}}}}
\newcommand{\KeywordTok}[1]{\textcolor[rgb]{0.13,0.29,0.53}{\textbf{#1}}}
\newcommand{\NormalTok}[1]{#1}
\newcommand{\OperatorTok}[1]{\textcolor[rgb]{0.81,0.36,0.00}{\textbf{#1}}}
\newcommand{\OtherTok}[1]{\textcolor[rgb]{0.56,0.35,0.01}{#1}}
\newcommand{\PreprocessorTok}[1]{\textcolor[rgb]{0.56,0.35,0.01}{\textit{#1}}}
\newcommand{\RegionMarkerTok}[1]{#1}
\newcommand{\SpecialCharTok}[1]{\textcolor[rgb]{0.81,0.36,0.00}{\textbf{#1}}}
\newcommand{\SpecialStringTok}[1]{\textcolor[rgb]{0.31,0.60,0.02}{#1}}
\newcommand{\StringTok}[1]{\textcolor[rgb]{0.31,0.60,0.02}{#1}}
\newcommand{\VariableTok}[1]{\textcolor[rgb]{0.00,0.00,0.00}{#1}}
\newcommand{\VerbatimStringTok}[1]{\textcolor[rgb]{0.31,0.60,0.02}{#1}}
\newcommand{\WarningTok}[1]{\textcolor[rgb]{0.56,0.35,0.01}{\textbf{\textit{#1}}}}
\usepackage{graphicx}
\makeatletter
\newsavebox\pandoc@box
\newcommand*\pandocbounded[1]{% scales image to fit in text height/width
  \sbox\pandoc@box{#1}%
  \Gscale@div\@tempa{\textheight}{\dimexpr\ht\pandoc@box+\dp\pandoc@box\relax}%
  \Gscale@div\@tempb{\linewidth}{\wd\pandoc@box}%
  \ifdim\@tempb\p@<\@tempa\p@\let\@tempa\@tempb\fi% select the smaller of both
  \ifdim\@tempa\p@<\p@\scalebox{\@tempa}{\usebox\pandoc@box}%
  \else\usebox{\pandoc@box}%
  \fi%
}
% Set default figure placement to htbp
\def\fps@figure{htbp}
\makeatother
\setlength{\emergencystretch}{3em} % prevent overfull lines
\providecommand{\tightlist}{%
  \setlength{\itemsep}{0pt}\setlength{\parskip}{0pt}}
\usepackage{amsmath}
\usepackage{amssymb}
\usepackage{booktabs}
\usepackage{graphicx}
\usepackage{float}
\usepackage{setspace}
\onehalfspacing
\usepackage{booktabs}
\usepackage{longtable}
\usepackage{array}
\usepackage{multirow}
\usepackage{wrapfig}
\usepackage{float}
\usepackage{colortbl}
\usepackage{pdflscape}
\usepackage{tabu}
\usepackage{threeparttable}
\usepackage{threeparttablex}
\usepackage[normalem]{ulem}
\usepackage{makecell}
\usepackage{xcolor}
\usepackage{bookmark}
\IfFileExists{xurl.sty}{\usepackage{xurl}}{} % add URL line breaks if available
\urlstyle{same}
\hypersetup{
  pdftitle={Taller 3 -- plan de muestreo estándar},
  pdfauthor={Diego Fernando Muñoz Portela 2415620, Valery Rivera Lopez 2417038, Alejandro Barros Rayo 2415837},
  hidelinks,
  pdfcreator={LaTeX via pandoc}}

\title{Taller 3 -- plan de muestreo estándar}
\author{Diego Fernando Muñoz Portela 2415620, Valery Rivera Lopez
2417038, Alejandro Barros Rayo 2415837}
\date{13 de marzo de 2026}

\begin{document}
\maketitle

{
\setcounter{tocdepth}{3}
\tableofcontents
}
\section{Punto 1}\label{punto-1}

\subsection{Enunciado}\label{enunciado}

El jefe de compras de una empresa ha exigido un AQL de 2.5\% para la
característica de un cierto producto que se recibe en lotes de 3000
unidades y se desea obtener un plan de inspección por muestreo usando
las normas Militar Standard.

\subsection{Preguntas}\label{preguntas}

\subsubsection{a. ¿Cuál es el plan de muestreo simple adecuado bajo los
niveles de inspección I, II, III, en los niveles normal, reducido y
severo?. Evalué la conveniencia de los planes resultantes. (OC, AOQ,
ATI).}\label{a.-cuuxe1l-es-el-plan-de-muestreo-simple-adecuado-bajo-los-niveles-de-inspecciuxf3n-i-ii-iii-en-los-niveles-normal-reducido-y-severo.-evaluuxe9-la-conveniencia-de-los-planes-resultantes.-oc-aoq-ati.}

Primero buuscamos cada uno de los \(n\),\(c\) y \(r\) para cada uno de
los planes simples (1,2,3) en cada uno de los niveles por medio de la
funcion:

La cual nos dara los siguientes resultados para el plan simple:

\begin{Shaded}
\begin{Highlighting}[]
\CommentTok{\# Crear el data frame sin columna Plan}
\NormalTok{tabla\_simple }\OtherTok{\textless{}{-}} \FunctionTok{data.frame}\NormalTok{(}
  \AttributeTok{Nivel =} \FunctionTok{c}\NormalTok{(}\StringTok{"I"}\NormalTok{, }\StringTok{"II"}\NormalTok{, }\StringTok{"III"}\NormalTok{,}
            \StringTok{"I"}\NormalTok{, }\StringTok{"II"}\NormalTok{, }\StringTok{"III"}\NormalTok{,}
            \StringTok{"I"}\NormalTok{, }\StringTok{"II"}\NormalTok{, }\StringTok{"III"}\NormalTok{),}
  \AttributeTok{n =} \FunctionTok{c}\NormalTok{(}\DecValTok{50}\NormalTok{, }\DecValTok{125}\NormalTok{, }\DecValTok{200}\NormalTok{, }\DecValTok{50}\NormalTok{, }\DecValTok{125}\NormalTok{, }\DecValTok{200}\NormalTok{, }\DecValTok{20}\NormalTok{, }\DecValTok{50}\NormalTok{, }\DecValTok{80}\NormalTok{),}
  \AttributeTok{c =} \FunctionTok{c}\NormalTok{(}\DecValTok{3}\NormalTok{, }\DecValTok{7}\NormalTok{, }\DecValTok{10}\NormalTok{, }\DecValTok{2}\NormalTok{, }\DecValTok{5}\NormalTok{, }\DecValTok{8}\NormalTok{, }\DecValTok{1}\NormalTok{, }\DecValTok{3}\NormalTok{, }\DecValTok{5}\NormalTok{),}
  \AttributeTok{r =} \FunctionTok{c}\NormalTok{(}\DecValTok{4}\NormalTok{, }\DecValTok{8}\NormalTok{, }\DecValTok{11}\NormalTok{, }\DecValTok{3}\NormalTok{, }\DecValTok{6}\NormalTok{, }\DecValTok{9}\NormalTok{, }\DecValTok{4}\NormalTok{, }\DecValTok{6}\NormalTok{, }\DecValTok{8}\NormalTok{)}
\NormalTok{)}

\FunctionTok{library}\NormalTok{(knitr)}
\FunctionTok{library}\NormalTok{(kableExtra)}

\NormalTok{tabla\_simple }\SpecialCharTok{\%\textgreater{}\%}
  \FunctionTok{kbl}\NormalTok{(}
    \AttributeTok{caption =} \StringTok{"Planes de muestreo simple según nivel de inspección"}\NormalTok{,}
    \AttributeTok{col.names =} \FunctionTok{c}\NormalTok{(}\StringTok{"Nivel"}\NormalTok{, }\StringTok{"n"}\NormalTok{, }\StringTok{"c"}\NormalTok{, }\StringTok{"r"}\NormalTok{),}
    \AttributeTok{align =} \StringTok{"c"}\NormalTok{,}
    \AttributeTok{booktabs =} \ConstantTok{TRUE}
\NormalTok{  ) }\SpecialCharTok{\%\textgreater{}\%}
  \FunctionTok{kable\_styling}\NormalTok{(}
  \AttributeTok{latex\_options =} \FunctionTok{c}\NormalTok{(}\StringTok{"striped"}\NormalTok{, }\StringTok{"hold\_position"}\NormalTok{),}
  \AttributeTok{position =} \StringTok{"center"}\NormalTok{,}
  \AttributeTok{font\_size =} \DecValTok{12}
\NormalTok{)}\SpecialCharTok{\%\textgreater{}\%}
  \FunctionTok{group\_rows}\NormalTok{(}\StringTok{"Normal"}\NormalTok{, }\DecValTok{1}\NormalTok{, }\DecValTok{3}\NormalTok{, }\AttributeTok{bold =} \ConstantTok{TRUE}\NormalTok{, }\AttributeTok{background =} \StringTok{"\#f0f0f0"}\NormalTok{) }\SpecialCharTok{\%\textgreater{}\%}
  \FunctionTok{group\_rows}\NormalTok{(}\StringTok{"Estricto"}\NormalTok{, }\DecValTok{4}\NormalTok{, }\DecValTok{6}\NormalTok{, }\AttributeTok{bold =} \ConstantTok{TRUE}\NormalTok{, }\AttributeTok{background =} \StringTok{"\#f0f0f0"}\NormalTok{) }\SpecialCharTok{\%\textgreater{}\%}
  \FunctionTok{group\_rows}\NormalTok{(}\StringTok{"Reducido"}\NormalTok{, }\DecValTok{7}\NormalTok{, }\DecValTok{9}\NormalTok{, }\AttributeTok{bold =} \ConstantTok{TRUE}\NormalTok{, }\AttributeTok{background =} \StringTok{"\#f0f0f0"}\NormalTok{) }\SpecialCharTok{\%\textgreater{}\%}
  \FunctionTok{column\_spec}\NormalTok{(}\DecValTok{1}\NormalTok{, }\AttributeTok{width =} \StringTok{"25\%"}\NormalTok{) }\SpecialCharTok{\%\textgreater{}\%}
  \FunctionTok{column\_spec}\NormalTok{(}\DecValTok{2}\NormalTok{, }\AttributeTok{width =} \StringTok{"25\%"}\NormalTok{) }\SpecialCharTok{\%\textgreater{}\%}
  \FunctionTok{column\_spec}\NormalTok{(}\DecValTok{3}\NormalTok{, }\AttributeTok{width =} \StringTok{"25\%"}\NormalTok{) }\SpecialCharTok{\%\textgreater{}\%}
  \FunctionTok{column\_spec}\NormalTok{(}\DecValTok{4}\NormalTok{, }\AttributeTok{width =} \StringTok{"25\%"}\NormalTok{) }\SpecialCharTok{\%\textgreater{}\%}
  \FunctionTok{add\_header\_above}\NormalTok{(}\FunctionTok{c}\NormalTok{(}\StringTok{" "} \OtherTok{=} \DecValTok{1}\NormalTok{, }\StringTok{"Parámetros del plan"} \OtherTok{=} \DecValTok{3}\NormalTok{)) }\SpecialCharTok{\%\textgreater{}\%}
  \FunctionTok{footnote}\NormalTok{(}\AttributeTok{general =} \StringTok{"n: tamaño de muestra; c: número de aceptación; r: número de rechazo"}\NormalTok{)}
\end{Highlighting}
\end{Shaded}

\begin{table}[!h]
\centering
\caption{\label{tab:tabla de c, n y r}Planes de muestreo simple según nivel de inspección}
\centering
\fontsize{12}{14}\selectfont
\begin{tabular}[t]{>{\centering\arraybackslash}p{25%}>{\centering\arraybackslash}p{25%}>{\centering\arraybackslash}p{25%}>{\centering\arraybackslash}p{25%}}
\toprule
\multicolumn{1}{c}{ } & \multicolumn{3}{c}{Parámetros del plan} \\
\cmidrule(l{3pt}r{3pt}){2-4}
Nivel & n & c & r\\
\midrule
\addlinespace[0.3em]
\multicolumn{4}{l}{\cellcolor[HTML]{f0f0f0}{\textbf{Normal}}}\\
\hspace{1em}\cellcolor{gray!10}{I} & \cellcolor{gray!10}{50} & \cellcolor{gray!10}{3} & \cellcolor{gray!10}{4}\\
\hspace{1em}II & 125 & 7 & 8\\
\hspace{1em}\cellcolor{gray!10}{III} & \cellcolor{gray!10}{200} & \cellcolor{gray!10}{10} & \cellcolor{gray!10}{11}\\
\addlinespace[0.3em]
\multicolumn{4}{l}{\cellcolor[HTML]{f0f0f0}{\textbf{Estricto}}}\\
\hspace{1em}I & 50 & 2 & 3\\
\hspace{1em}\cellcolor{gray!10}{II} & \cellcolor{gray!10}{125} & \cellcolor{gray!10}{5} & \cellcolor{gray!10}{6}\\
\hspace{1em}III & 200 & 8 & 9\\
\addlinespace[0.3em]
\multicolumn{4}{l}{\cellcolor[HTML]{f0f0f0}{\textbf{Reducido}}}\\
\hspace{1em}\cellcolor{gray!10}{I} & \cellcolor{gray!10}{20} & \cellcolor{gray!10}{1} & \cellcolor{gray!10}{4}\\
\hspace{1em}II & 50 & 3 & 6\\
\hspace{1em}\cellcolor{gray!10}{III} & \cellcolor{gray!10}{80} & \cellcolor{gray!10}{5} & \cellcolor{gray!10}{8}\\
\bottomrule
\multicolumn{4}{l}{\rule{0pt}{1em}\textit{Note: }}\\
\multicolumn{4}{l}{\rule{0pt}{1em}n: tamaño de muestra; c: número de aceptación; r: número de rechazo}\\
\end{tabular}
\end{table}

opcion 2

\begin{Shaded}
\begin{Highlighting}[]
\CommentTok{\# Crear el data frame original, pero cambiando la columna Plan a I, II, III}
\NormalTok{tabla\_simple }\OtherTok{\textless{}{-}} \FunctionTok{data.frame}\NormalTok{(}
  \AttributeTok{Tipo =} \FunctionTok{c}\NormalTok{(}\StringTok{"Normal"}\NormalTok{, }\StringTok{"Normal"}\NormalTok{, }\StringTok{"Normal"}\NormalTok{,}
           \StringTok{"Estricto"}\NormalTok{, }\StringTok{"Estricto"}\NormalTok{, }\StringTok{"Estricto"}\NormalTok{,}
           \StringTok{"Reducido"}\NormalTok{, }\StringTok{"Reducido"}\NormalTok{, }\StringTok{"Reducido"}\NormalTok{),}
  \AttributeTok{Plan =} \FunctionTok{factor}\NormalTok{(}\FunctionTok{c}\NormalTok{(}\StringTok{"I"}\NormalTok{, }\StringTok{"II"}\NormalTok{, }\StringTok{"III"}\NormalTok{, }\StringTok{"I"}\NormalTok{, }\StringTok{"II"}\NormalTok{, }\StringTok{"III"}\NormalTok{, }\StringTok{"I"}\NormalTok{, }\StringTok{"II"}\NormalTok{, }\StringTok{"III"}\NormalTok{),}
                \AttributeTok{levels =} \FunctionTok{c}\NormalTok{(}\StringTok{"I"}\NormalTok{, }\StringTok{"II"}\NormalTok{, }\StringTok{"III"}\NormalTok{)),}
  \AttributeTok{n =} \FunctionTok{c}\NormalTok{(}\DecValTok{50}\NormalTok{, }\DecValTok{125}\NormalTok{, }\DecValTok{200}\NormalTok{, }\DecValTok{50}\NormalTok{, }\DecValTok{125}\NormalTok{, }\DecValTok{200}\NormalTok{, }\DecValTok{20}\NormalTok{, }\DecValTok{50}\NormalTok{, }\DecValTok{80}\NormalTok{),}
  \AttributeTok{c =} \FunctionTok{c}\NormalTok{(}\DecValTok{3}\NormalTok{, }\DecValTok{7}\NormalTok{, }\DecValTok{10}\NormalTok{, }\DecValTok{2}\NormalTok{, }\DecValTok{5}\NormalTok{, }\DecValTok{8}\NormalTok{, }\DecValTok{1}\NormalTok{, }\DecValTok{3}\NormalTok{, }\DecValTok{5}\NormalTok{),}
  \AttributeTok{r =} \FunctionTok{c}\NormalTok{(}\DecValTok{4}\NormalTok{, }\DecValTok{8}\NormalTok{, }\DecValTok{11}\NormalTok{, }\DecValTok{3}\NormalTok{, }\DecValTok{6}\NormalTok{, }\DecValTok{9}\NormalTok{, }\DecValTok{4}\NormalTok{, }\DecValTok{6}\NormalTok{, }\DecValTok{8}\NormalTok{)}
\NormalTok{)}

\FunctionTok{library}\NormalTok{(knitr)}
\FunctionTok{library}\NormalTok{(kableExtra)}

\NormalTok{tabla\_simple }\SpecialCharTok{\%\textgreater{}\%}
  \FunctionTok{kbl}\NormalTok{(}
    \AttributeTok{caption =} \StringTok{"Planes de muestreo simple según nivel de inspección"}\NormalTok{,}
    \AttributeTok{col.names =} \FunctionTok{c}\NormalTok{(}\StringTok{"Nivel"}\NormalTok{, }\StringTok{"Nivel del plan"}\NormalTok{, }\StringTok{"n"}\NormalTok{, }\StringTok{"c"}\NormalTok{, }\StringTok{"r"}\NormalTok{),}
    \AttributeTok{align =} \StringTok{"c"}\NormalTok{,}
    \AttributeTok{booktabs =} \ConstantTok{TRUE}
\NormalTok{  ) }\SpecialCharTok{\%\textgreater{}\%}
  \FunctionTok{kable\_styling}\NormalTok{(}
  \AttributeTok{latex\_options =} \FunctionTok{c}\NormalTok{(}\StringTok{"striped"}\NormalTok{, }\StringTok{"hold\_position"}\NormalTok{),}
  \AttributeTok{position =} \StringTok{"center"}\NormalTok{,}
  \AttributeTok{font\_size =} \DecValTok{12}
\NormalTok{  ) }\SpecialCharTok{\%\textgreater{}\%}
  \CommentTok{\# Agrupar por Tipo (Normal, Estricto, Reducido)}
  \FunctionTok{group\_rows}\NormalTok{(}\StringTok{"Normal"}\NormalTok{, }\DecValTok{1}\NormalTok{, }\DecValTok{3}\NormalTok{, }\AttributeTok{bold =} \ConstantTok{TRUE}\NormalTok{, }\AttributeTok{background =} \StringTok{"\#f0f0f0"}\NormalTok{) }\SpecialCharTok{\%\textgreater{}\%}
  \FunctionTok{group\_rows}\NormalTok{(}\StringTok{"Estricto"}\NormalTok{, }\DecValTok{4}\NormalTok{, }\DecValTok{6}\NormalTok{, }\AttributeTok{bold =} \ConstantTok{TRUE}\NormalTok{, }\AttributeTok{background =} \StringTok{"\#f0f0f0"}\NormalTok{) }\SpecialCharTok{\%\textgreater{}\%}
  \FunctionTok{group\_rows}\NormalTok{(}\StringTok{"Reducido"}\NormalTok{, }\DecValTok{7}\NormalTok{, }\DecValTok{9}\NormalTok{, }\AttributeTok{bold =} \ConstantTok{TRUE}\NormalTok{, }\AttributeTok{background =} \StringTok{"\#f0f0f0"}\NormalTok{) }\SpecialCharTok{\%\textgreater{}\%}
  \CommentTok{\# Resaltar la columna del nivel del plan}
  \FunctionTok{column\_spec}\NormalTok{(}\DecValTok{2}\NormalTok{, }\AttributeTok{bold =} \ConstantTok{TRUE}\NormalTok{) }\SpecialCharTok{\%\textgreater{}\%}
  \CommentTok{\# Añadir cabecera superior para los parámetros}
  \FunctionTok{add\_header\_above}\NormalTok{(}\FunctionTok{c}\NormalTok{(}\StringTok{" "} \OtherTok{=} \DecValTok{2}\NormalTok{, }\StringTok{"Parámetros del plan"} \OtherTok{=} \DecValTok{3}\NormalTok{)) }\SpecialCharTok{\%\textgreater{}\%}
  \FunctionTok{footnote}\NormalTok{(}\AttributeTok{general =} \StringTok{"n: tamaño de muestra; c: número de aceptación; r: número de rechazo"}\NormalTok{)}
\end{Highlighting}
\end{Shaded}

\begin{table}[!h]
\centering
\caption{\label{tab:deded}Planes de muestreo simple según nivel de inspección}
\centering
\fontsize{12}{14}\selectfont
\begin{tabular}[t]{c>{}cccc}
\toprule
\multicolumn{2}{c}{ } & \multicolumn{3}{c}{Parámetros del plan} \\
\cmidrule(l{3pt}r{3pt}){3-5}
Nivel & Nivel del plan & n & c & r\\
\midrule
\addlinespace[0.3em]
\multicolumn{5}{l}{\cellcolor[HTML]{f0f0f0}{\textbf{Normal}}}\\
\hspace{1em}\cellcolor{gray!10}{Normal} & \textbf{\cellcolor{gray!10}{I}} & \cellcolor{gray!10}{50} & \cellcolor{gray!10}{3} & \cellcolor{gray!10}{4}\\
\hspace{1em}Normal & \textbf{II} & 125 & 7 & 8\\
\hspace{1em}\cellcolor{gray!10}{Normal} & \textbf{\cellcolor{gray!10}{III}} & \cellcolor{gray!10}{200} & \cellcolor{gray!10}{10} & \cellcolor{gray!10}{11}\\
\addlinespace[0.3em]
\multicolumn{5}{l}{\cellcolor[HTML]{f0f0f0}{\textbf{Estricto}}}\\
\hspace{1em}Estricto & \textbf{I} & 50 & 2 & 3\\
\hspace{1em}\cellcolor{gray!10}{Estricto} & \textbf{\cellcolor{gray!10}{II}} & \cellcolor{gray!10}{125} & \cellcolor{gray!10}{5} & \cellcolor{gray!10}{6}\\
\hspace{1em}Estricto & \textbf{III} & 200 & 8 & 9\\
\addlinespace[0.3em]
\multicolumn{5}{l}{\cellcolor[HTML]{f0f0f0}{\textbf{Reducido}}}\\
\hspace{1em}\cellcolor{gray!10}{Reducido} & \textbf{\cellcolor{gray!10}{I}} & \cellcolor{gray!10}{20} & \cellcolor{gray!10}{1} & \cellcolor{gray!10}{4}\\
\hspace{1em}Reducido & \textbf{II} & 50 & 3 & 6\\
\hspace{1em}\cellcolor{gray!10}{Reducido} & \textbf{\cellcolor{gray!10}{III}} & \cellcolor{gray!10}{80} & \cellcolor{gray!10}{5} & \cellcolor{gray!10}{8}\\
\bottomrule
\multicolumn{5}{l}{\rule{0pt}{1em}\textit{Note: }}\\
\multicolumn{5}{l}{\rule{0pt}{1em}n: tamaño de muestra; c: número de aceptación; r: número de rechazo}\\
\end{tabular}
\end{table}

\begin{Shaded}
\begin{Highlighting}[]
\CommentTok{\# Crear data frame con los valores definitivos}
\NormalTok{tabla\_doble\_final }\OtherTok{\textless{}{-}} \FunctionTok{data.frame}\NormalTok{(}
  \AttributeTok{Tipo =} \FunctionTok{rep}\NormalTok{(}\FunctionTok{c}\NormalTok{(}\StringTok{"Normal"}\NormalTok{, }\StringTok{"Estricto"}\NormalTok{, }\StringTok{"Reducido"}\NormalTok{), }\AttributeTok{each =} \DecValTok{6}\NormalTok{),}
  \AttributeTok{Plan =} \FunctionTok{rep}\NormalTok{(}\FunctionTok{rep}\NormalTok{(}\DecValTok{1}\SpecialCharTok{:}\DecValTok{3}\NormalTok{, }\AttributeTok{each =} \DecValTok{2}\NormalTok{), }\AttributeTok{times =} \DecValTok{3}\NormalTok{),}
  \AttributeTok{Etapa =} \FunctionTok{rep}\NormalTok{(}\DecValTok{1}\SpecialCharTok{:}\DecValTok{2}\NormalTok{, }\AttributeTok{times =} \DecValTok{9}\NormalTok{),}
  \AttributeTok{n =} \FunctionTok{c}\NormalTok{(}\DecValTok{32}\NormalTok{, }\DecValTok{32}\NormalTok{, }\DecValTok{80}\NormalTok{, }\DecValTok{80}\NormalTok{, }\DecValTok{125}\NormalTok{, }\DecValTok{125}\NormalTok{,   }\CommentTok{\# Normal}
        \DecValTok{32}\NormalTok{, }\DecValTok{32}\NormalTok{, }\DecValTok{80}\NormalTok{, }\DecValTok{80}\NormalTok{, }\DecValTok{125}\NormalTok{, }\DecValTok{125}\NormalTok{,   }\CommentTok{\# Estricto}
        \DecValTok{13}\NormalTok{, }\DecValTok{13}\NormalTok{, }\DecValTok{32}\NormalTok{, }\DecValTok{32}\NormalTok{, }\DecValTok{50}\NormalTok{, }\DecValTok{50}\NormalTok{),    }\CommentTok{\# Reducido}
  \AttributeTok{c =} \FunctionTok{c}\NormalTok{(}\DecValTok{1}\NormalTok{, }\DecValTok{4}\NormalTok{, }\DecValTok{3}\NormalTok{, }\DecValTok{8}\NormalTok{, }\DecValTok{5}\NormalTok{, }\DecValTok{12}\NormalTok{,          }\CommentTok{\# Normal}
        \DecValTok{0}\NormalTok{, }\DecValTok{3}\NormalTok{, }\DecValTok{2}\NormalTok{, }\DecValTok{6}\NormalTok{, }\DecValTok{3}\NormalTok{, }\DecValTok{11}\NormalTok{,          }\CommentTok{\# Estricto}
        \DecValTok{0}\NormalTok{, }\DecValTok{1}\NormalTok{, }\DecValTok{1}\NormalTok{, }\DecValTok{4}\NormalTok{, }\DecValTok{2}\NormalTok{, }\DecValTok{6}\NormalTok{),          }\CommentTok{\# Reducido}
  \AttributeTok{r =} \FunctionTok{c}\NormalTok{(}\DecValTok{4}\NormalTok{, }\DecValTok{5}\NormalTok{, }\DecValTok{7}\NormalTok{, }\DecValTok{9}\NormalTok{, }\DecValTok{9}\NormalTok{, }\DecValTok{13}\NormalTok{,          }\CommentTok{\# Normal}
        \DecValTok{3}\NormalTok{, }\DecValTok{4}\NormalTok{, }\DecValTok{5}\NormalTok{, }\DecValTok{7}\NormalTok{, }\DecValTok{7}\NormalTok{, }\DecValTok{12}\NormalTok{,          }\CommentTok{\# Estricto}
        \DecValTok{4}\NormalTok{, }\DecValTok{5}\NormalTok{, }\DecValTok{4}\NormalTok{, }\DecValTok{7}\NormalTok{, }\DecValTok{6}\NormalTok{, }\DecValTok{9}\NormalTok{)           }\CommentTok{\# Reducido}
\NormalTok{)}

\CommentTok{\# Mostrar el data frame}
\FunctionTok{print}\NormalTok{(tabla\_doble\_final)}
\end{Highlighting}
\end{Shaded}

\begin{verbatim}
##        Tipo Plan Etapa   n  c  r
## 1    Normal    1     1  32  1  4
## 2    Normal    1     2  32  4  5
## 3    Normal    2     1  80  3  7
## 4    Normal    2     2  80  8  9
## 5    Normal    3     1 125  5  9
## 6    Normal    3     2 125 12 13
## 7  Estricto    1     1  32  0  3
## 8  Estricto    1     2  32  3  4
## 9  Estricto    2     1  80  2  5
## 10 Estricto    2     2  80  6  7
## 11 Estricto    3     1 125  3  7
## 12 Estricto    3     2 125 11 12
## 13 Reducido    1     1  13  0  4
## 14 Reducido    1     2  13  1  5
## 15 Reducido    2     1  32  1  4
## 16 Reducido    2     2  32  4  7
## 17 Reducido    3     1  50  2  6
## 18 Reducido    3     2  50  6  9
\end{verbatim}

\end{document}
